
\chapter{Einrichtung des Arduino UNO}


In diesem Abschnitt wird der Vorgang beschrieben, wie der Arduino UNO
mit einem Computer verbunden und in der Arduino IDE eingerichtet wird.
Dazu werden die notwendigen Pakete installiert, das Board sowie der
zugehörige Port ausgewählt und abschließend anhand eines
Beispiel-Sketches überprüft, ob die Einrichtung erfolgreich war.

\section{Installation der Pakete für den Arduino UNO}

Um den Arduino UNO zu programmieren, wird zunächst das passende
Treiberpaket in der Arduino IDE installiert. Dazu wird in der Menüleiste
der Punkt Tools ausgewählt und darunter der Eintrag Board Manager
geöffnet. Im Suchfeld des Board Managers wird "Arduino UNO" eingegeben
\ref{pic:ARV}. Aus den Suchergebnissen wird das Paket "Arduino
AVR Boards" ausgewählt und über die Schaltfläche "Install"
installiert.

\includegraphics[width=1\textwidth,keepaspectratio]{Arduino/Setup/AVR}
\captionof{figure}{Installation der Arduino AVR Boards}
\label{pic:ARV}



Das Paket Arduino AVR Boards enthält die Unterstützung für mehrere
Boards der AVR-Familie, unter anderem für den Arduino UNO. Nach
Abschluss der Installation wird der Arduino UNO über ein
USB-A-auf-USB-B-Kabel mit dem Computer verbunden.

\section{Verbindung und Auswahl des Boards}

Nachdem der Arduino UNO über das USB-Kabel mit dem Computer verbunden
wurde, wird die Arduino IDE geöffnet. Es erscheint eine leere
Umgebungsseite mit den standardmäßigen Funktionen void setup() und void
loop(). Anschließend wird in der Menüleiste unter dem Punkt Tools der
Unterpunkt Board ausgewählt und dort der Arduino UNO als angeschlossenes
Board festgelegt (vgl. Abbildung 1.25).

Abbildung 1.25: Auswahl des angeschlossenen Boards -- hier Arduino UNO

Mit dieser Auswahl ist die Arduino IDE darauf vorbereitet, den Arduino
UNO zu erkennen und mit ihm zu kommunizieren.

\section{Auswahl des Ports}

Nach der Boardauswahl wird der verbundene Port überprüft. Im Gegensatz
zu einigen anderen Arduino-Boards muss der Arduino UNO dafür nicht
zusätzlich in einen Bootloader-Modus versetzt werden, der Port wird beim
Anschließen automatisch erkannt. Dazu wird in der Menüleiste unter dem
Punkt Tools der Unterpunkt Port ausgewählt und der verfügbare Port für
das Hochladen des Programms festgelegt (vgl. Abbildung 1.26).

Abbildung 1.26: Auswahl des verfügbaren Ports zum Hochladen des Arduino-Sketches
\section{Hochladen und Überprüfen eines Sketches}

Nachdem der passende Port ausgewählt wurde, kann das Arduino-Programm
hochgeladen werden. Vor dem Hochladen wird der Code zunächst überprüft
(verifiziert), um sicherzustellen, dass keine Fehler oder Warnungen
vorliegen. Anschließend wird der Code über die Schaltfläche Upload auf
das Board übertragen (vgl. Abbildung 1.27). Ob der Arduino UNO korrekt
mit der Arduino IDE verbunden ist, lässt sich anhand der Meldung in der
unteren rechten Ecke des IDE-Fensters überprüfen, welche das verbundene
Board und den Port anzeigt.

 Abbildung 1.27: Hochladen des Programms auf das Arduino Board

Nach dem Hochladen wird der Code kompiliert und etwaige Fehler werden im
unteren Ausgabefenster angezeigt. Sobald der Code erfolgreich übertragen
wurde, wird der Port erneut über Tools und Port überprüft (vgl.
Abbildung 1.28), um die Ausgabe im Serial Monitor anzeigen zu können.


Abbildung 1.28: Einstellen des Ports

\section{Test der Konfiguration anhand des Beispiels Blink}

Um zu überprüfen, ob das Board korrekt mit der Arduino IDE verbunden
ist, wird das eingebaute Beispiel \emph{Blink} verwendet. Dieser
einfache Test lässt die auf dem Board integrierte LED blinken und
bestätigt damit sowohl die Verbindung als auch die grundlegende
Funktionalität des Arduino UNO.

Dazu wird in der Arduino IDE über das Menü File der Punkt Examples
ausgewählt. Unter den eingebauten Beispielen wird unter 01.Basics der
Eintrag Blink ausgewählt \ref{pic:Blink}. Das Beispiel-Sketch
öffnet sich daraufhin in einem neuen Fenster innerhalb der IDE. Über die Schaltfläche Upload in der Werkzeugleiste wird das Beispiel kompiliert und auf den Arduino UNO übertragen. Nach
Abschluss des Uploads beginnt die integrierte LED zu blinken.

\includegraphics[width=1\textwidth,keepaspectratio]{Arduino/Setup/Blink}
\captionof{figure}{Beispiel-Sketch Blink}
\label{pic:Blink}

\captionof{code}{Blink Programm \label{CodeIDE}}
\ArduinoExternal{}{../../Code/Blink/Blink.ino}

Mit dem erfolgreichen Blinken der integrierten LED ist bestätigt, dass
der Arduino UNO ordnungsgemäß mit der Arduino IDE verbunden ist und
seine grundlegende Funktionalität überprüft wurde. Das Board ist damit
bereit für die weitere Entwicklung.

